\addchap{Preface}
\label{ch:about}

This is \fillin{one paragraph: what the book is (a reading guide to the papers that explain X; a crash course in
Y), the question it answers, and who it is written for}.

\section*{What it gives you}

\begin{enumerate}[label=(\arabic*)]
  \item A \emph{curated, tiered} list of the important papers, each with a short card saying what it does, why it
        matters and what to read it for (Part~\ref{part:literature}).
  \item \fillin{one item per part: what the reader gets from it, with \textbackslash ref to the part}.
  \item The \emph{papers themselves}. Part~\ref{part:papers} lists the arXiv papers cited in the book (\npapers{} in
        all), with links; in the combined edition the papers follow that list, joined into the same file, and the
        bookmarks lead to each of them.
\end{enumerate}

\section*{How the list was chosen and checked}

\textbf{Selection.} \fillin{where the candidates came from (the compiler's knowledge of the field, the most widely
used surveys, citation counts retrieved from Semantic Scholar on <date>) and the criteria: influence (did others
build on it?), pivot (did it change a design choice?), clarity, recency}.

\textbf{Verification.} \fillin{the title, authors, dates and category of every arXiv paper were fetched from the
arXiv API on <date>; venues come from arXiv comment fields, publisher records (Crossref) or proceedings pages;
every other source was checked against a publisher, library or official page; and how the claims were checked:
name only checks that were actually done, e.g.\ independent readers against the papers, with reports in
provenance/}.

\textbf{Caveats.} \fillin{(1) numbers are as reported by the authors and were not reproduced; (2) preprints may
change; (3) the tiers and the narrative are the compiler's judgment; (4) anything else a careful reader should
know}.

\section*{Source files}

The LaTeX sources and the scripts that generate the paper cards and the bibliography are kept with this PDF, and
the directory \texttt{provenance/} holds \fillin{the evidence that exists: reading notes, writers' reports,
fact-check reports and the log of applied fixes; list only what is there}.

The look of the book (fonts, colors, boxes, running heads, the structure of the front matter) follows that of
Evan Chen's \emph{An Infinitely Large Napkin}~\cite{chen2024napkin}; the book uses none of its text or artwork.

\addchap{Advice for the reader}
\label{ch:advice}
\setcounter{section}{0}
\renewcommand{\thesection}{\arabic{section}}

\section{Prerequisites}

The book assumes \fillin{what the reader must already know}. It does not assume \fillin{what it teaches from
zero}; \cref{ch:tour} \fillin{supplies what the papers assume}.

\section{Deciding what to read}

There is no need to read this book in order.

\begin{description}[leftmargin=!,labelwidth=4.7cm,font=\sffamily\bfseries\color{maroon},itemsep=2pt]
  \item[Part~\ref{part:start}: Getting started] \fillin{one line per part, naming its chapters with \textbackslash cref}.
  \item[Part~\ref{part:literature}: The literature] \fillin{one line}.
  \item[Part~\ref{part:papers}: The papers] The papers cited in the book, in the order listed in that part.
\end{description}

The chart below shows how the chapters depend on each other. Arrows mean ``read first'', dotted lines are optional
or partial dependencies, and the chapters with thick borders form the short path of \S\ref{sec:paths}.
\textbf{Pick a chapter you find interesting, and follow the arrows backward to see what to read first.}

% Chapter dependencies, in the style of the chart in the Napkin's "Advice for the reader".
% Solid arrow: read first. Dotted arrow: optional or partial dependency. Thick box: the short reading path.
% One node per chapter (or per block of a lettered part); nodes link to their chapters.
\begin{center}
\begingroup
\hypersetup{linkcolor=black}
\begin{tikzpicture}[x=1cm,y=0.86cm,
  cn/.style={draw=black,fill=yellow!20!white,rounded corners=2.5pt,line width=0.5pt,inner sep=2.5pt,
             minimum width=3.45cm,minimum height=0.92cm,align=center,font=\small},
  key/.style={line width=1.5pt},
  dep/.style={-{Stealth[length=2.1mm,width=1.6mm]},line width=0.45pt},
  opt/.style={-{Stealth[length=2.1mm,width=1.6mm]},line width=0.55pt,dash pattern=on 0.3pt off 2.2pt,line cap=round},
  hdr/.style={font=\sffamily\bfseries\footnotesize,text=crimson,anchor=south}]
  \node[hdr] at (2.0,1.25) {Part~\ref{part:start}};
  \node[hdr] at (7.4,1.25) {Part~\ref{part:literature}};
  \node[cn,key] (c1) at (2.0,0) {{\scriptsize\textbf{Ch \ref*{ch:tour}}}\\\hyperref[ch:tour]{Exponential decay}};
  \node[cn,key] (c2) at (7.4,0) {{\scriptsize\textbf{Ch \ref*{ch:reading}}}\\\hyperref[ch:reading]{Reading the literature}};
  \node[cn]     (g)  at (12.8,0) {{\scriptsize\textbf{App.\ \ref*{app:glossary}}}\\\hyperref[app:glossary]{Glossary}};
  \draw[dep] (c1) -- (c2);
  \draw[opt] (c2) -- (g);
\end{tikzpicture}
\endgroup
\end{center}


\section{Reading tiers}

Every paper with a card carries a tier, shown by stars and by the kind of box that holds the card:

\begin{center}
\small
\begin{tabularx}{\textwidth}{@{}l l l X@{}}
\swatchA & \textcolor{crimson}{\ding{72}\ding{72}\ding{72}} & \sffamily\bfseries\textcolor{tierA}{Essential} & Read it (or at least its key sections) if you read only a handful. \\[3pt]
\swatchB & \textcolor{crimson}{\ding{72}\ding{72}} & \sffamily\bfseries\textcolor{tierB}{Important} & Introduces an idea that later work builds on; read when that thread matters to you. \\[3pt]
\swatchC & \textcolor{crimson}{\ding{72}} & \sffamily\bfseries\textcolor{tierC}{Context} & Useful background or a notable variant; skim or look up. \\
\end{tabularx}
\end{center}

The tiers are judgments, informed by citation counts (as a rough signal of influence) and by whether a paper
changed what others did. They are the compiler's opinion, and they will age.

\section{Reading paths}\label{sec:paths}

\paragraph{The short path.} \fillin{the few chapters and cards to read if the reader has an hour or two}.

\paragraph{The core path.} \fillin{one sentence: how to use the table}. One row per paper in teaching order, at most
twelve.

\begin{center}
\small\renewcommand{\arraystretch}{1.3}
\begin{tabularx}{\textwidth}{@{}>{\sffamily\bfseries\color{crimson}}r@{\hspace{1.1em}}>{\raggedright\arraybackslash}p{5.3cm}X@{}}
\toprule
\# & Paper & Read it for \\
\midrule
1 & \fillin{name~\textbackslash cite\{key\}} & \fillin{the idea it introduces} \\
2 & \fillin{name~\textbackslash cite\{key\}} & \fillin{the idea it introduces} \\
\bottomrule
\end{tabularx}
\end{center}

\section{Conventions and notations}

\fillin{the notation used throughout, defined once}.

Boxes are numbered within each section, as in a textbook: Paper~2.1.1 is the first box of Section~2.1. Citations in
brackets (for example [Che24]) are labels of the entries in the bibliography at the end: the first letters of the
authors' names, a plus sign if there are more than three authors, and the year. The bibliography lists, after each
entry, the pages that cite it.

\section{Further reading}

\fillin{where to keep up with the field: surveys, courses, code, datasets, venues}.
